\chapter{The editor}
\label{ch:editor}

\section{The toolkit}
\code{editor/ui.*} is an immediate-mode toolkit (\code{ui::Context}): buttons, fields with expressions (\code{+=1},
\code{*=2}, \code{L(0,10)}, \code{R(0,1)}), dropdowns, colour fields and the Color window, popups, menus, tooltips,
tab strips, toasts and the job tray, drawn by the canvas (\code{render/canvas.*}: anti-aliased TrueType text from the
OS fonts; icons rasterized once per size). Every widget records its rectangle, which is what the Fitts audit reads and
what the coverage map's \code{--widgets} counts. Text contrast is measured per theme against WCAG's 4.5:1
\cite{W3C2018Web}; polygroup and signal colours follow Okabe and Ito's colour-universal set \cite{Okabe2008Color}.

\section{Docking, windows and the module registry}
\code{editor/editor.cpp} holds a dock tree of areas and tabs: drag a tab onto a tab bar to group it, onto an edge to
split, out of the editor for a window of its own on any display (\adr{0059}); a layout is a dock string plus a focus
(the tool groups a stage shows, \adr{0027}), and the user's layouts and the whole arrangement go into a
\code{.blendity-ui} file. The module registry (\code{editor/modules.cpp}, \adr{0073}) is one row per window kind: title,
icon, group, the stages it belongs to, whether a second window may hold one, whether its group folds, its shortcut.
The 26 windows are Scene, Game, Hierarchy, Inspector, Project, Console, Learn, Research (retired), Profiler, Render, UV
Editor, Materials, Modeling Tools, Lighting, Textures, Rigging, Animation, Sequencer, Live Link, Editing, My Tools,
Draw, Layers, History, Board and Avatar; the tab strips, the Window menu, the layouts and the tests read the registry,
so a new window is one row.

\section{Keymaps, feels and the modal pattern}
\code{editor/keymap.cpp}: \code{shortcut\_actions()} lists every action (id, label, category, context);
\code{keymap\_preset(name)} builds the bindings for Unity, Unreal, Blender, Maya, 3ds Max, SketchUp and Adobe;
\code{set\_editor\_feel} sets the preset and the Scene view's mouse navigation together; the user's overrides sit on top;
\code{handle\_shortcuts} matches a key against the most specific active context first (Scene + Edit, Edit, Scene,
anywhere). Every shortcut label in a menu or tooltip comes from the keymap through \code{shortcut\_text}, so no label can
lie under another preset (\adr{0085}). Every edit operator follows one modal pattern (\code{docs/USABILITY.md}): a key
starts it, the mouse or a typed number shapes it, X / Y / Z lock an axis, Ctrl snaps, a click or Enter confirms, Esc
cancels leaving no undo step, F9 adjusts afterwards as the same step. The Wheel (\adr{0068}) is one radial menu for
Space, the pen's button or a finger held still, with a set per mode the user fills from any menu; My Tools panels and
Assign Shortcut give any command a button or a key (\adr{0069}).

\section{Tools}
The Scene view runs the active tool each frame: the Draw tools with SketchUp inference and guides (\code{draw\_tool.cpp}; the Materials window's preview spheres are
Blinn--Phong \cite{Phong1975}),
the Knife, Push/Pull as a modal operator, G / R / S (\code{modal\_transform.cpp}), the gizmo with Shift-extrude and
origin editing, the sketch pen with its guides, canvases and the camera solved from the perspective guide
(\adr{0063}--\adr{0067}, \cite{Igarashi1997Interactive, Chen2013Sweep, Gryaditskaya2020Lifting}), Sculpt Mode
(\adr{0074}--\adr{0078}, \cite{Jacobson2013Robust, Gibson1998Constrained, Huang2018Quadriflow}), Paint Mode by
projection (\adr{0081}, \cite{Baxter2004Impasto}), Weight Paint, Strokes and Scatter (\adr{0084}), Puppet Posing's
handles with a learned pose prior and exact IK (\adr{0087}, \cite{Oreshkin2022Protores}), and the VR View's
controllers through OpenXR (\adr{0096}, \cite{OpenXR2026}). Fingers and pens come through one pointer model
(\adr{0062}): in Tablet Mode fingers navigate and the pen works; on Linux through XInput~2; on a Mac the trackpad.

\section{The console and headless runs}
Every tool is reachable from the Console's command line (\code{editor\_commands.cpp}; \code{help} lists about 240
commands), and that is the editor's scripting surface: the tests, \code{--headless-screenshot --cmd}, the HCI
production runs, \code{tools/unity\_compare.py} and \code{unity\_e2e.py} all drive it. The console recalls its
history, answers an unknown command with the nearest known ones, and since \rnd{310} no verb an artist uses lives only
there: a test keeps every one reachable from a menu, a button or a field.

\section{The Learn tab}
\code{editor/learn.cpp} renders 45 lessons from \code{lessons.cpp}, each connecting a Unity concept, the Blender
concept with the source folder it lives in, and the theory in the reference books \cite{MarschnerShirley2021,
Gregory2018vol1, Gregory2018vol2}, with Try-it buttons that run console commands, questions that reveal their
answers, and a References lesson of the papers behind the features. Every new feature adds a lesson; the user guide
names the lesson for each chapter, and Index C of that guide lists each lesson's sources.

\section{Where the HCI work lives}
\adr{0085} set the methods (Nielsen's heuristics \cite{Nielsen1994Enhancing, Nielsen1993}, Fitts's law
\cite{Fitts1954Information}, the keystroke-level model \cite{Card1980Keystroke, Card1983}, cognitive walkthroughs
\cite{Polson1992Cognitive}, transfer of learning, WCAG) and the foundations every area uses: toasts with a Fix and
disabled buttons that say why, one job tray, status hints from a table, labels from the keymap, unsaved-changes
prompts, a History window, a colour-vision-safe palette and a High Contrast theme. The measuring sticks are in the
test suite (the Fitts dump, the coverage map, the label audit, the hint completeness test, the contrast test) and the
artefacts in \code{docs/hci/} (walkthroughs, the KLM table recounted by \code{tools/hci\_klm.py}, the comparison reports
against Blender, Unity, Godot, Unreal, Adobe and VRChat's limits, the 118-task checklist, the study protocol written
but not yet run \cite{Brooke1996Sus}).
